\documentclass[11pt]{article}
\usepackage[margin=1in]{geometry}
\usepackage{amsmath,amssymb,amsthm,mathtools}
\usepackage{array,booktabs,longtable}
\usepackage[hidelinks]{hyperref}

\newtheorem{theorem}{Theorem}[section]
\newtheorem{lemma}[theorem]{Lemma}
\newtheorem{proposition}[theorem]{Proposition}
\newtheorem{corollary}[theorem]{Corollary}
\theoremstyle{remark}
\newtheorem{remark}[theorem]{Remark}
\newcommand{\splus}{s^+}
\newcommand{\sminus}{s^-}
\newcommand{\one}{\mathbf 1}
\newcommand{\kappad}{\kappa}

\title{Positive Square Energy for Rank-Seven Kernels of Order Eight}
\author{AI-generated manuscript; no human author is claimed}
\date{11 November 2026}

\begin{document}
\maketitle

\begin{abstract}
Let $G$ be a finite simple connected graph of cyclomatic rank seven whose
block decomposition has exactly one positive-rank cyclic block.  Suppose that,
after suppressing the degree-two vertices of that block, its loopless
$2$-connected multigraph kernel has eight vertices.  We prove
$\splus(G)\geq |V(G)|$, where $\splus$ is the sum of the squares of the
positive adjacency eigenvalues.  Recent general theorems already cover every
tree-free member of this class; the additional scope here is the uniform
allowance of arbitrary rooted-tree attachments, including leaf-bearing
instances.  The proof combines an exact census of the $4{,}015$ order-eight
kernels and their $26{,}426{,}026$ physical parity orbits with analytic coarse
certificates and a disjoint exact certificate partition of the $493{,}417$
residual orbits.  The last residual is handled by an induced partition into two
tricyclic territories.  The executable theorem gate pins every proof source
and artifact and offers a quick integrity mode and a full read-only replay
mode.  Three expensive historical base-producer outputs remain explicit
immutable trust roots; full mode freshly exact-replays the 29,000
direct-rational witness records rather than retaining them as a fourth root.
No claim is made for kernel orders nine through twelve, for multiblock
rank-seven graphs, or for all connected heptacyclic graphs.
\end{abstract}

\section{Statement, scope, and spectral tools}

All graphs in this paper are finite, simple, and undirected.  If the adjacency
eigenvalues of a graph $H$ are $\lambda_1,\ldots,\lambda_n$, write
\[
 \splus(H)=\sum_{\lambda_i>0}\lambda_i^2,
 \qquad
 \sminus(H)=\sum_{\lambda_i<0}\lambda_i^2.
\]
A connected graph has cyclomatic rank
$\beta(H)=|E(H)|-|V(H)|+1$.  A \emph{cyclic block} means a maximal
$2$-connected block of positive cyclomatic rank; bridges are rank-zero blocks.

\begin{theorem}[Scoped order-eight kernel theorem]\label{thm:main}
Let $G$ be a finite simple connected graph with $\beta(G)=7$.  Assume that
$G$ has exactly one positive-rank cyclic block $B$.  Suppress every degree-two
vertex of $B$, and let $K$ be the resulting loopless $2$-connected multigraph.
If $|V(K)|=8$, then
\[
 \boxed{\splus(G)\geq |V(G)|}.
\]
Equivalently, the conclusion holds for every simple subdivision of every
rank-seven order-eight kernel, with arbitrary finite rooted trees attached at
branch vertices or subdivision vertices.
\end{theorem}

The hypotheses in Theorem~\ref{thm:main} are indispensable to the scope of
this manuscript.  They do not include rank-seven graphs with two or more
positive-rank blocks, and they do not include kernels of orders nine through
twelve.  In particular, the theorem is not a theorem about all connected
graphs with $|E|=|V|+6$.

\paragraph{Relation to the square-energy literature.}
Akbari, Kumar, Mohar, Pragada, and Zhang conjectured that every connected graph
of order $n$ and size $m\geq n+1$ satisfies $\splus(G)\geq n$
\cite{AKMPZ}.  A connected graph of cyclomatic rank seven has $m=n+6$, so
Theorem~\ref{thm:main} verifies only the restricted class stated above.
Akbari, Hu, and Liu subsequently proved the conclusion for every
$2$-connected noncycle \cite{AHLTwoConnected} and, more generally, for every
connected noncycle of minimum degree at least two \cite{AHLMinDegree}.
Consequently every tree-free member of our class is already covered by those
results.  Within the one-cyclic-block class considered here, any nontrivial
finite rooted-tree attachment has a non-root leaf, so the resulting graph has
minimum degree one and is outside the latter theorem.  The certificate
argument below treats all such attachments uniformly.  The general
$n-1$ square-energy theorem \cite{LTZ} gives a baseline one unit below the
claim here.  We make no claim of firstness: the dated search audit
\cite{NoveltyAudit} located no previous theorem with the full combined scope,
but a negative search does not establish novelty or priority.

We use two standard spectral devices.  The first is the induced partition
inequality established in the square- and $p$-energy literature
\cite{LinearLowerBound,VertexPartitioning}: positive square energy is
superadditive over induced vertex territories.

\begin{lemma}[Induced superadditivity]\label{lem:super}
If $V(H)=V_1\mathbin{\dot\cup}\cdots\mathbin{\dot\cup}V_r$, then
\[
 \splus(H)\geq\sum_{j=1}^r\splus(H[V_j]).
\]
\end{lemma}

\begin{proof}
For a real symmetric matrix $A$,
\[
 \operatorname{tr}(A_+^2)
 =\max_{X\succeq0}\bigl(2\operatorname{tr}(AX)-\operatorname{tr}(X^2)\bigr).
\]
For $A=A(H)$, test this variational formula with the block-diagonal matrix
whose blocks are $A(H[V_j])_+$.  Cross-territory entries of $A$ do not
contribute, and the displayed inequality follows.
\end{proof}

Second, for an edge-containing graph $H$, define the DNN constant
\[
 \kappad(H)=
 \min_{\substack{R\succeq0\\R_{vv}=1}}
 \sum_{uv\in E(H)}\frac{2}{1-R_{uv}},                       \tag{1}
\]
with boundary values interpreted by limits.  The correlation dual for the
optimized doubly nonnegative bound gives
\begin{equation}
 \sminus(H)\leq\kappad(H).                                  \tag{2}
\end{equation}
Moreover, $\kappad$ is additive under one-vertex sums and
$\kappad(T)=|E(T)|$ for every tree $T$; proofs are recorded in the companion
tricyclic manuscript~\cite{Tricyclic}.  These statements may also be obtained
from the primal DNN formulation by orthogonally gluing the Gram data at a cut
vertex.

\section{Suppression and the finite frontier}

Let $B$ be a $2$-connected graph of cyclomatic rank seven.  Suppressing its
degree-two vertices gives a loopless $2$-connected multigraph $K$ of minimum
degree at least three, and
\[
 |E(K)|=|V(K)|+6,
 \qquad
 3|V(K)|\leq2|E(K)|.
\]
Consequently $2\leq |V(K)|\leq12$.  Conversely, $B$ is obtained by replacing
each physical edge of $K$ by a positive-length path, with pairwise disjoint
interiors.  Simplicity permits at most one unit path in each parallel class.
Every part of $G$ outside $B$ is a rooted tree attached at one vertex of $B$.
The complete removable-ear generation theorem and its exact audit give
$4{,}015$ order-eight kernels~\cite{CensusTheorem}.

Fix a rank-seven kernel $K$ of order eight.  It has fourteen physical edges.
For a parallel class of multiplicity $m$ containing $o$ odd replacement
paths, define its canonical simple lengths, up to permutation within the
class, by
\[
 o=0:\ (2,\ldots,2),
 \qquad
 o>0:\ (1,3,\ldots,3,2,\ldots,2).                          \tag{3}
\]
Thus every simple realization in the same physical parity orbit can be
ordered as $\ell=c+2d$ with $d\geq0$ coordinatewise.  For a canonical vector
$c=(c_1,\ldots,c_{14})$, put
\[
 \mathcal F(c)=\{c\}\cup\{c+2e_i:1\leq i\leq14\}.          \tag{4}
\]

The exact path atom explains both the DNN budget and the all-length lift.

\begin{lemma}[Exact path elimination]\label{lem:path}
Fix unit branch vectors with correlation $r$ at the endpoints of a path of
length $q\geq1$.  Its least excess over the baseline contribution $q$ in
(1) is
\[
 f_q(r)=q\tan^2\!\left(
 \frac{\arccos((-1)^q r)}{2q}\right).                       \tag{5}
\]
For fixed $r$ and fixed parity,
$f_{q+2}(r)\leq f_q(r)$, with strict inequality unless the transformed
endpoint angle vanishes.
\end{lemma}

\begin{proof}
Alternately negate the vectors along the path.  The transformed endpoints
have angular distance $\beta=\arccos((-1)^q r)$.  The spherical triangle
inequality and convexity of $\tan^2(x/2)$ show that $q$ equal angular steps
minimize the sum, yielding (5).  For real $q>0$, differentiation of
$q\tan^2(\beta/(2q))$ has the sign of
$\sin z\cos z-2z$, where $z=\beta/(2q)$; this is negative for $z>0$.
Endpoint cases follow by limits.
\end{proof}

Suppose a retained branch Gram has total path excess at most six.  If its
subdivision core has $L$ edges, then $\kappad(B)\leq L+6$.  If $t$ tree edges
are attached, one-vertex additivity gives
$\kappad(G)\leq L+6+t$.  Since rank seven implies $|V(B)|=L-6$, equations
(2) and $\splus(G)+\sminus(G)=2|E(G)|$ give
\begin{equation}
 \splus(G)\geq2(L+t)-(L+6+t)=L-6+t=|V(G)|.                  \tag{6}
\end{equation}
Lemma~\ref{lem:path} promotes a canonical or coordinate-frontier Gram to every
same-parity realization that dominates it.  A structural owner is admitted
only when it supplies its own all-length and rooted-tree argument.

\section{Exact census and owner partition}

The order-eight census enumerates all physical parity rows and quotients them
by the full multigraph automorphism group.  It independently checks every
materialized orbit count with Burnside's lemma.  The resulting totals are
\begin{center}
\begin{tabular}{lr}
\toprule
quantity & exact count\\
\midrule
order-eight kernels & $4{,}015$\\
physical parity rows & $33{,}006{,}092$\\
automorphism orbits & $26{,}426{,}026$\\
coarse-certified orbits & $25{,}932{,}609$\\
coarse residual orbits & $493{,}417$\\
residual targets in (4) & $7{,}401{,}255$\\
\bottomrule
\end{tabular}
\end{center}
The exact manifest and its eight authenticated chunks satisfy
$25{,}932{,}609+493{,}417=26{,}426{,}026$.  The coarse certificate is a
four-color regular-tetrahedron Gram whose exact scaled cost is at most the
rank-seven excess budget six; hence its accepted rows are already DNN owners.
The remaining $493{,}417$ rows are partitioned by the theorem gate as follows.

\begin{table}[ht]
\centering
\small
\begin{tabular}{lrr}
\toprule
exclusive owner & residual orbits & targets\\
\midrule
payload-free symbolic lanes & $605$ & $9{,}075$\\
direct rational witness packs & $29{,}000$ & $435{,}000$\\
scalar SOS, after precedence & $194{,}351$ & $2{,}915{,}265$\\
typed-diagonal Gram, after precedence & $185{,}309$ & $2{,}779{,}635$\\
boundary shared-Gram rescue & $296$ & $4{,}440$\\
structural cycle-Gram lanes I--III & $245$ & $3{,}675$\\
packet lanes IV--XII & $29{,}013$ & $435{,}195$\\
global post-XII shared packets & $54{,}597$ & $818{,}955$\\
final induced-territory owner & $1$ & $15$\\
\midrule
exact total & $493{,}417$ & $7{,}401{,}255$\\
\bottomrule
\end{tabular}
\caption{Disjoint precedence partition of the coarse residual frontier.}
\label{tab:owners}
\end{table}

The row identity in Table~\ref{tab:owners} is
\begin{equation}\tag{7}
\begin{split}
493{,}417={}&605+29{,}000+194{,}351+185{,}309+296\\
            &+245+29{,}013+54{,}597+1.
\end{split}
\end{equation}
The first five lanes leave $83{,}856$ rows.  Structural lanes I--III leave
$83{,}611$.  Packet lanes IV--XII leave $54{,}598$, and the global post-XII
lane leaves one.  Thus the exact remainder chain used by the master gate is
\[
 83{,}856\longrightarrow83{,}611\longrightarrow54{,}598
 \longrightarrow1.                                         \tag{8}
\]

We summarize why the entries are theorem owners rather than merely successful
numerical searches.

\begin{itemize}
\item The payload-free lanes use rank-one signings, diagonally dominant signed
imbalance Grams, or exact simplex/mixed-atom PSD completions.
\item The scalar and typed-diagonal lanes express every accepted Gram as an
exact sum of rational Gram squares and a nonnegative diagonal completion.
Twenty gap-free typed receipts cover all $492{,}812$ rows after the
payload-free lane; the precedence bridge verifies the exact exclusive sets.
\item Every direct, rescue, structural, and packet witness is accepted only
after exact rational replay.  For a shared packet, numerical optimization is
proposal-only; acceptance verifies the unit branch Gram and each path-chain
cost with rational arithmetic.
\item The structural verifier replays all $245$ stored owners and all
$3{,}675$ frontier costs.  The packet-chain verifier replays the nine exact
complements IV--XII, whose owner counts sum to $29{,}013$.
\item Each DNN or packet owner retains a branch Gram and invokes
Lemma~\ref{lem:path}, so every same-parity length vector is covered.  New tree
vertices can be assigned the negative of their parent vector, adding one to
both $\kappad$ and the edge count.  This gives arbitrary rooted-tree
attachments.
\end{itemize}

The immutable trust boundary is important.  Three expensive historical
producer searches remain roots: (1) payload-free recognizer semantics, (2) raw
scalar-SOS membership semantics, and (3) the $492{,}812$ typed searches behind
the twenty receipts.  The base precedence bridge reports the historical
four-root boundary, but the master discharges its fourth item by freshly
canonical-decoding and exact-\texttt{Fraction} replaying every one of the
$29{,}000$ direct-rational witness records, covering $435{,}000$ targets.  The
full master also reconstructs the selected $492{,}812$-tuple projection in
memory from the authenticated census chunks and checks the pinned cache
field-for-field; it neither writes a newly regenerated cache artifact nor
reruns the three searches just named.  Separate read-only verifiers replay the
structural and packet witnesses.  The final claim is therefore a complete
exact certificate chain relative to three explicitly named and byte-pinned
base producer outputs, not a claim that every historical search was rerun.

\section{The unique residual orbit}

After the post-XII replay, the unique residual is the six-field row
\[
 [416536,416987,6832,E,p,2],                                  \tag{9}
\]
where
\begin{align*}
 E={}&\{03,06,07,12,14,15,24,25,36,37,45,47,56,67\},\\
 p={}&(1,1,1,1,1,1,1,1,1,1,1,1,0,1).
\end{align*}
The first two entries are respectively its stream index and residual source
index, $6832$ is the global kernel index, and the orbit size is two.  Its
support consists of two vertex-disjoint copies of $K_4$, on
\[
 A=\{0,3,6,7\},\qquad C=\{1,2,4,5\},
\]
joined by the two physical paths $P_{74}$ and $P_{65}$.  The following
proposition is stronger than the fifteen finite obligations in (4).

\begin{proposition}[Two tricyclic territories]\label{prop:singleton}
Start with the two displayed copies of $K_4$, replace all twelve internal
$K_4$ edges and the two connectors $74,65$ by arbitrary positive-length
paths with pairwise disjoint interiors, and attach arbitrary finite rooted
trees at branch or subdivision vertices.  The resulting graph $H$ satisfies
$\splus(H)\geq|V(H)|$, with no parity restriction on any path.
\end{proposition}

\begin{proof}
Let $C_A$ and $C_C$ be the two $K_4$ subdivisions.  Orient both connector
paths from the $A$ side to the $C$ side.  Put in $U_0$ all vertices of $C_A$
and all internal connector vertices, and put in $W_0$ all vertices of $C_C$.
Every attached rooted tree follows its root, producing a vertex partition
$V(H)=U\mathbin{\dot\cup}W$.

The induced graph $H[U]$ consists of $C_A$, the two connector stubs obtained by
deleting their terminal $C$-side branch vertices, and rooted forests meeting
this cyclic core in one root.  The graph $H[W]$ consists of $C_C$ and its
assigned rooted forests.  Both induced graphs are connected.  This description
also covers a unit connector: then its sole edge crosses the partition and it
has no internal stub.

If $C_A$ has $L_A$ edges and its retained forest has $t_A$ edges, then a
subdivision of $K_4$ has $L_A-2$ vertices, whence
\[
 |E(H[U])|=L_A+t_A=|V(H[U])|+2.
\]
The identical calculation holds for $H[W]$.  The tricyclic theorem says that
every finite simple connected graph $J$ with
$|E(J)|=|V(J)|+2$ satisfies $\splus(J)\geq|V(J)|$
\cite{Tricyclic}.  Therefore Lemma~\ref{lem:super} gives
\[
 \splus(H)\geq\splus(H[U])+\splus(H[W])
 \geq|U|+|W|=|V(H)|.
\]
No edge-addition monotonicity, connector debit, or DNN gluing is used.
\end{proof}

The singleton artifact binds (9) byte-for-byte to
Proposition~\ref{prop:singleton}.  Because the proposition covers all lengths,
all parities, and arbitrary rooted trees, it supplies the last owner in (7).

\section{Proof of the theorem}

\begin{proof}[Proof of Theorem~\ref{thm:main}]
Suppress $B$ to its rank-seven order-eight kernel $K$.  The exhaustive kernel
census and physical automorphism quotient place every parity state of $K$ in
exactly one of two sets.

If it is coarse-certified, its tetrahedral correlation Gram has excess at most
six and equation (6) applies.  Otherwise it is one of the $493{,}417$ residual
orbits.  The disjoint exact identity (7) assigns it to exactly one row of
Table~\ref{tab:owners}.  Every owner except the last supplies an exact DNN or
packet certificate together with the all-length and rooted-tree lift described
above, so again (6) applies.  The unique last orbit is covered uniformly by
Proposition~\ref{prop:singleton}.  These alternatives exhaust every kernel,
physical parity orbit, simple subdivision length, and rooted-tree attachment
within the stated scope.  Hence $\splus(G)\geq|V(G)|$.
\end{proof}

\section{Executable verification and integrity boundary}

The authoritative integration gate is
\begin{verbatim}
research/rank-seven-order-eight-kernel-theorem-verifier.py
\end{verbatim}
with SHA-256
\begin{verbatim}
8e2baed4455f0721a1f670bef2f5a97d937fc0baa2f5f5f88ee38d6709565255
\end{verbatim}
on the verified integration bytes.  The repository preserves the complete
source lineage in its merge history.  From the repository root, run
\begin{verbatim}
python3 -B research/rank-seven-order-eight-kernel-theorem-verifier.py
python3 -O -B research/rank-seven-order-eight-kernel-theorem-verifier.py
python3 -B research/rank-seven-order-eight-kernel-theorem-verifier.py \
  --print-manifest
\end{verbatim}
The quick text-only output has SHA-256
\begin{verbatim}
193981b59dde97ba0dfd680e2ffd7477f02663c0645fb024ffa120a38fbb9afe
\end{verbatim}
and the quick manifest-plus-text output has SHA-256
\begin{verbatim}
b5922982b4c14ef9a80ed5c71bbb01cbfdc887daf94bd662774b9339415de041.
\end{verbatim}
The gate executes its own optimized master child, requires byte-identical
normal and \texttt{-O} output, rejects 21 focused hostile mutations, and checks
that 26 proof-critical Python sources contain no optimization-removable
\texttt{assert}.

Quick mode is an integrity and composition gate: it hashes every child source
and mathematical dependency, authenticates the post-XII compressed and raw
streams, validates strict schemas and the singleton bytes, and pins receipts
from prior exact normal/\texttt{-O} child runs.  It deliberately does not
rerun the expensive semantic children.  To do so, run
\begin{verbatim}
python3 -B research/rank-seven-order-eight-kernel-theorem-verifier.py --full
\end{verbatim}
Full mode successfully executes the base bridge, structural replay,
packet-chain replay, post-XII replay, exhaustive kernel/census regeneration,
all 29,000 direct-rational witness replays, singleton reconstruction, and the
imported tricyclic finite gates both normally and under \texttt{-O}.  It
reconstructs the selected 492,812-row census projection in memory and validates
the pinned cache field-for-field, but does not write a newly regenerated cache
artifact or rerun the three residual historical producer searches.  The
post-XII normalized proof receipt is pinned as
\begin{verbatim}
1dbd2337043e283e4b8913b89151da6b000cbc516a822f3b5f24bbf94e89eb4b.
\end{verbatim}
The distinction between quick integrity verification, full read-only replay,
and fresh producer recomputation is part of the theorem contract.

Build this manuscript with
\begin{verbatim}
bash scripts/build-paper.sh rank-seven-order-eight-kernels
\end{verbatim}

\paragraph{Scoped nonclaim.}
Theorem~\ref{thm:main} concerns only cyclomatic-rank-seven graphs with exactly
one positive-rank cyclic block and with suppressed kernel order eight.  It
makes no claim for kernel orders nine through twelve, for any multiblock
rank-seven graph, or for all connected heptacyclic graphs.  Its tree-free
subfamily was already covered by \cite{AHLTwoConnected,AHLMinDegree}; this
paper's additional scope relative to those results is the uniform treatment of
arbitrary rooted-tree attachments within the stated finite kernel frontier.
It makes no firstness or unconditional novelty claim.  It does not claim fresh
reruns of the three residual historical base producer searches.  It uses
neither parity-changing subdivision monotonicity nor edge-addition
monotonicity.

\section*{AI disclosure}
This manuscript and its proof-object synthesis were generated by an AI coding
and mathematical reasoning system from the cited repository proofs and exact
certificate artifacts.  No human authorship, human review, or human
verification is claimed.  The result should be assessed from the displayed
arguments, explicit trust boundary, frozen sources, and reproducible checks.
No external publication was performed as part of this manuscript build.

\begin{thebibliography}{99}
\bibitem{AKMPZ}
S.~Akbari, H.~Kumar, B.~Mohar, S.~Pragada, and S.~Zhang,
\emph{Refinement of a conjecture on positive square energy of graphs},
arXiv:2506.07264v1, 2025.

\bibitem{AHLTwoConnected}
S.~Akbari, F.-T.~Hu, and Y.-Y.~Liu,
\emph{Positive and negative square energies of $2$-connected graphs},
arXiv:2609.04069v1, 2026.

\bibitem{AHLMinDegree}
S.~Akbari, F.-T.~Hu, and Y.-Y.~Liu,
\emph{Positive square energy of graphs with minimum degree at least two},
arXiv:2609.06778v1, 2026.

\bibitem{LTZ}
Y.~Liu, Q.~Tang, and S.~Zhang,
\emph{The positive and negative square-energy conjecture},
arXiv:2607.18031v1, 2026.

\bibitem{VertexPartitioning}
S.~Akbari, H.~Kumar, B.~Mohar, and S.~Pragada,
\emph{Vertex partitioning and $p$-energy of graphs},
arXiv:2503.16882v2, 2025.

\bibitem{LinearLowerBound}
S.~Akbari, H.~Kumar, B.~Mohar, and S.~Pragada,
\emph{A linear lower bound for the square energy of graphs},
Electron. J. Combin. \textbf{32}(3) (2025), Paper P3.53;
doi:10.37236/13467.

\bibitem{EFGW}
C.~Elphick, M.~Farber, F.~Goldberg, and P.~Wocjan,
\emph{Conjectured bounds for the sum of squares of positive eigenvalues of a
graph}, Discrete Math. \textbf{339} (2016), 2215--2223;
doi:10.1016/j.disc.2016.01.021.

\bibitem{AbiadEtAl}
A.~Abiad, L.~de Lima, D.~N.~Desai, K.~Guo, L.~Hogben, and J.~Madrid,
\emph{Positive and negative square energies of graphs},
Electron. J. Linear Algebra \textbf{39} (2023), 307--326.

\bibitem{NoveltyAudit}
\emph{Rank-seven order-eight kernel source and status audit}, 11 November 2026;
\path{research/rank-seven-order-eight-source-status-audit-2026-11-11.md}.

\subsection*{Proof artifacts and reproducibility}

\bibitem{Tricyclic}
Companion manuscript, \emph{Positive Square Energy of Every Tricyclic Graph},
2026; \path{all-tricyclic-graphs/paper.tex}.  The final gate pins the
manuscript, PDF, README, hostile audit, and finite dependency verifiers by
SHA-256.

\bibitem{CensusTheorem}
\emph{Rank-seven removable-ear census theorem and exact audit}, 2026;
\path{positive-square-energy/heptacyclic-general/rank-seven-removable-ear-census-theorem.md}
and \path{research/rank-seven-kernel-frontier-census.py}.

\bibitem{Reduction}
\emph{Rank-seven single-block structural theorem: exact reduction and
obstructions}, 2026;
\path{positive-square-energy/heptacyclic-general/rank-seven-single-block-structural-theorem.md}.

\bibitem{BaseBridge}
Exact base-owner precedence bridge and hostile audit,
\path{positive-square-energy/experiments/rank7_order8_typed_precedence_bridge.py}.

\bibitem{StructuralReplay}
Exact read-only replay for structural lanes I--III,
\path{positive-square-energy/experiments/rank7_order8_structural_exact_verifier.py}.

\bibitem{PacketReplay}
Exact read-only replay for packet lanes IV--XII,
\path{positive-square-energy/experiments/verify_rank7_order8_packet_chain.py}.

\bibitem{Singleton}
\emph{Final order-eight residual: two induced tricyclic territories}, 2026;
\path{positive-square-energy/heptacyclic-general/rank-seven-order-eight-final-two-k4-territory.md}.
\end{thebibliography}

\end{document}
